\documentclass[11pt,a4paper]{article}
\usepackage[margin=25mm]{geometry}
\usepackage{amsmath,amssymb}
\usepackage{booktabs}
\usepackage{microtype}
\setlength{\parindent}{0pt}
\setlength{\parskip}{0.6em}
\newcommand{\Prob}{\mathbb{P}}

\title{All-On Allocation Model with Overlapping Intervals}
\author{}
\date{}

\begin{document}
\maketitle

\section{Circuit}
A circuit has \(N\) customers, each of class \(j\in\{S,M,L\}\) with current
\(I_j\) (A) and load \(P_j=V I_j\) (kW) at voltage \(V\). A composition
\(\boldsymbol n=(n_S,n_M,n_L)\), \(n_S+n_M+n_L=N\), occurs with probability
\begin{equation}
  w(\boldsymbol n)=\frac{N!}{n_S!\,n_M!\,n_L!}\,p_S^{n_S}p_M^{n_M}p_L^{n_L},
  \qquad p_S+p_M+p_L=1 .
\end{equation}

\section{Allocation}
There are \(K\in\{1,2,3\}\) free intervals \([s_k,e_k)\), \(k=1,\dots,K\),
which may overlap. Each medium and large customer is assigned exactly one
interval. Let \(m_k\) and \(l_k\) be the medium and large customers in
interval \(k\):
\begin{equation}
  \sum_{k=1}^K m_k=n_M,\qquad \sum_{k=1}^K l_k=n_L,\qquad m_k,l_k\in\mathbb{Z}_{\ge0}.
\end{equation}
Inside its interval a customer draws its own class current and load; outside
it, it behaves as a small customer. Small customers never change.

\section{Segments}
Only time inside at least one interval is analysed; time outside every
interval is ignored. Cut this time at every interval edge. Each resulting
segment \(g=[a_g,b_g)\) is covered by a non-empty set of intervals
\(C(g)=\{k: s_k\le a_g,\ b_g\le e_k\}\); overlaps give \(|C(g)|\ge2\).
During segment \(g\),
\begin{equation}
  M_g=\sum_{k\in C(g)} m_k,\qquad
  L_g=\sum_{k\in C(g)} l_k,\qquad
  S_g=N-M_g-L_g .
\end{equation}
The circuit current and the energy consumed in the segment are
\begin{align}
  A_g &= I_S S_g + I_M M_g + I_L L_g \quad\text{(A)},\\
  E_g &= \bigl(P_S S_g + P_M M_g + P_L L_g\bigr)\,(b_g-a_g) \quad\text{(kWh)}.
\end{align}

\section{Solar soak}
Let \(X_d(t)\) be the VIC1 utility-solar curtailment (MW) on day \(d\), slot \(t\)
(\(\Delta t=5\) min). If a share \(u\) of the \(N_{\mathrm{VIC}}\) customers join the
tariff, \(C=uN_{\mathrm{VIC}}/N\) circuits share it, so one circuit has
\begin{equation}
  x_d(t)=X_d(t)/C .
\end{equation}
Curtailment is what remains after existing demand, so only load added by the
free intervals can soak it. In slot \(t\in g\) this is
\begin{equation}
  D(t)=(P_M-P_S)\,M_g+(P_L-P_S)\,L_g \quad\text{(kW)}.
\end{equation}
The solar soak per circuit per day, over the days \(\mathcal D\) (one chosen day, or
all 365 days of 2025), is
\begin{equation}
  \text{soak}=\frac{1}{|\mathcal D|}\sum_{d\in\mathcal D}\sum_t \min\bigl(D(t),\,x_d(t)\bigr)\,\Delta t
  \quad\text{(kWh/day)}.
\end{equation}
The minimum is taken day by day before averaging; since \(\min\) is concave,
using the mean curtailment profile instead would overstate soak.

Hourly outputs of the single-circuit run sum a kW series \(f\) over the slots in each
bin \(B\): \(\sum_{t\in B} f(t)\,\Delta t\), with \(f\) the consumed load
\(N P_S + D(t)\), the added load \(D(t)\), the curtailment share \(x_d(t)\), or the
soak \(\min(D,x_d)\), each averaged over \(\mathcal D\).

\section{Overload}
A circuit with composition \(\boldsymbol n\) and allocation \((\boldsymbol m,\boldsymbol l)\)
overloads if the current exceeds the limit \(T\) in any segment:
\begin{equation}
  r(\boldsymbol n,\boldsymbol m,\boldsymbol l)=
  \begin{cases}1, & \max_g A_g > T,\\ 0, & \text{otherwise.}\end{cases}
\end{equation}

\section{Auto allocation and violation probability}
The auto allocation minimises the peak current,
\begin{equation}
  (\boldsymbol m^*,\boldsymbol l^*)(\boldsymbol n)
  =\arg\min_{\boldsymbol m,\boldsymbol l}\ \max_g A_g ,
\end{equation}
checking every integer split (ties: the first split found). If every circuit uses its
auto allocation, the probability that a circuit overloads is
\begin{equation}
  \Prob(\text{violation})=\sum_{\boldsymbol n} w(\boldsymbol n)\,
  r\bigl(\boldsymbol n,\boldsymbol m^*(\boldsymbol n),\boldsymbol l^*(\boldsymbol n)\bigr),
\end{equation}
skipping compositions with \(w(\boldsymbol n)<10^{-15}\). A composition has a safe
allocation if and only if its auto allocation is safe.

\section{Pareto front}
Both objectives are averages over compositions. Each circuit may use its own
allocation \(a(\boldsymbol n)\):
\begin{align}
  O_{\mathrm{safe}} &= 1-\sum_{\boldsymbol n} w(\boldsymbol n)\, r\bigl(\boldsymbol n,a(\boldsymbol n)\bigr)
      && \text{(share of circuits without overload)},\\
  O_{\mathrm{soak}} &= \sum_{\boldsymbol n} w(\boldsymbol n)\, \text{soak}\bigl(\boldsymbol n,a(\boldsymbol n)\bigr)
      && \text{(expected soak per circuit, kWh/day)},
\end{align}
with soak averaged over all days of 2025. Both are to be maximised.

For each composition every integer allocation is checked, keeping
\begin{align}
  \sigma_{\max}(\boldsymbol n) &= \max_a \text{soak}(\boldsymbol n,a), &
  \sigma_{\mathrm{safe}}(\boldsymbol n) &= \max_{a:\,r(\boldsymbol n,a)=0} \text{soak}(\boldsymbol n,a).
\end{align}
Because \(r\in\{0,1\}\), only these two allocations can lie on the front. If
several allocations reach \(\sigma_{\max}\), a safe one is preferred.
Start with every composition on its \(\sigma_{\max}\) allocation (most soak).
Composition \(\boldsymbol n\) can then be moved to its safe allocation, which raises
\(O_{\mathrm{safe}}\) by \(w(\boldsymbol n)\) and lowers \(O_{\mathrm{soak}}\) by
\(w(\boldsymbol n)\,\Delta(\boldsymbol n)\), where
\begin{equation}
  \Delta(\boldsymbol n)=\sigma_{\max}(\boldsymbol n)-\sigma_{\mathrm{safe}}(\boldsymbol n).
\end{equation}
Moving compositions in increasing order of \(\Delta(\boldsymbol n)\) traces the
Pareto front; the segment between two consecutive points is reached by moving
only part of the circuits with that composition. The front ends at
\(O_{\mathrm{safe}}=1-\Prob(\text{violation})\) of the auto allocation.

Along the front the share of medium and large customers placed in interval \(k\) is
\begin{equation}
  \frac{\sum_{\boldsymbol n} w(\boldsymbol n)\,\bigl(m_k+l_k\bigr)\bigl(a(\boldsymbol n)\bigr)}
       {\sum_{\boldsymbol n} w(\boldsymbol n)\,(n_M+n_L)} .
\end{equation}

With one interval there is a single allocation, so the front is one point.
Every allocation with fewer intervals is also available with more when the
interval sets are nested, so the front for \(K+1\) intervals is never below
the front for \(K\).

The uptake \(u\) sets \(x_d(t)=X_d(t)/C\): at low uptake each circuit has more
curtailment than added load, so soak is limited by load; at high uptake
curtailment is the limit and spreading load over more hours soaks more.

\section{Computation}
Soak depends on an allocation only through \((M_g,L_g)\) in each segment, so for every
segment a table is precomputed once:
\begin{align}
  F_g(M,L)&=\frac{1}{|\mathcal D|}\sum_{d}\sum_{t\in g}
    \min\bigl((P_M-P_S)M+(P_L-P_S)L,\ x_d(t)\bigr)\,\Delta t,\\
  \text{soak}(\boldsymbol n,a)&=\sum_g F_g(M_g,L_g).
\end{align}
Each composition has
\(\binom{n_M+K-1}{K-1}\binom{n_L+K-1}{K-1}\) allocations; all are evaluated
(\(9.4\times10^6\) over all compositions for \(K=3\)).

\begin{center}
\small
\begin{tabular}{ll}
\toprule
Quantity & Code (\texttt{play/})\\
\midrule
Inputs, \(P_j=VI_j\), \(C\) & \texttt{config.py}: \texttt{load\_kw}, \texttt{circuits}\\
\(x_d(t)\) & \texttt{solar.py}: \texttt{excess\_per\_circuit}\\
Segments, \(S_g,M_g,L_g\), \(A_g\), \(w(\boldsymbol n)\) & \texttt{circuit.py}: \texttt{segments}, \texttt{active}, \texttt{peak}, \texttt{weight}\\
Auto allocation, \(\Prob(\text{violation})\) & \texttt{circuit.py}: \texttt{auto\_allocation}, \texttt{violation\_probability}\\
\(D(t)\), hourly energy, \(F_g\) & \texttt{soak.py}: \texttt{load\_profiles}, \texttt{kwh}, \texttt{soak\_table}\\
\(\sigma_{\max},\sigma_{\mathrm{safe}}\), front & \texttt{frontier.py}: \texttt{evaluate}, \texttt{pareto\_front}\\
Single-circuit run, Pareto run & \texttt{run\_allocation.py}, \texttt{run\_pareto.py}\\
\bottomrule
\end{tabular}
\end{center}

\section{Inputs}
Shared inputs are in \texttt{play/config.py}; \(K\), \(\boldsymbol n\), the allocation and the
day are set per run in \texttt{play/run\_allocation.py}, and the uptakes compared in
\texttt{play/run\_pareto.py}.
\begin{center}
\begin{tabular}{lll}
\toprule
Symbol & Code variable & Default\\
\midrule
\(N\) & \texttt{N} & 40\\
\(T\) & \texttt{T} & 1000 A\\
\(I_S,I_M,I_L\) & \texttt{CURRENT} & 20, 40, 100 A\\
\(V\) & \texttt{V} & 240 V\\
\(u\) & \texttt{UPTAKE} & 0.1\\
\(N_{\mathrm{VIC}}\) & \texttt{VIC\_CUSTOMERS} & 2{,}900{,}000\\
\(\mathcal D\) & \texttt{SOLAR\_DAY} & 2025-11-06 (Pareto: all of 2025)\\
\(p_S,p_M,p_L\) & \texttt{P\_TYPE} & 0.69, 0.30, 0.01\\
\([s_k,e_k)\) & \texttt{INTERVAL\_SETS} & 11--14 \,/\, 8--11, 11--14 \,/\, 8--11, 10:30--13:30, 11--14\\
\(K\) & \texttt{NUM\_INTERVALS} & 3\\
\(\boldsymbol n\) & \texttt{COMPOSITION} & (32, 7, 1)\\
\((\boldsymbol m,\boldsymbol l)\) & \texttt{ALLOCATION} & auto\\
\bottomrule
\end{tabular}
\end{center}

\end{document}
